\QHeading{第1题　口算。}
\par\noindent


\begin{oralgrid}
\OralItem{1}{2.3+0.6}{2.9}
\OralItem{2}{5.3+3.2}{8.5}
\OralItem{3}{7.5-3.5}{4}
\OralItem{4}{1.4-0.5}{0.9}
\OralItem{5}{1-0.2}{0.8}
\OralItem{6}{3.5-3.1}{0.4}
\OralItem{7}{0.54+1.44}{1.98}
\OralItem{8}{0.2+1.8}{2}
\OralItem{9}{8.8-8.7}{0.1}
\OralItem{10}{3.26+5.41}{8.67}
\OralItem{11}{7.58+2.35}{9.93}
\OralItem{12}{9.87-4.52}{5.35}\end{oralgrid}
\par

\QHeading{口算思路补充（对应上表题号）}

\Answer{上面列出得数，下面按小题给出不同计算思路。计数单位法适用于小数，整数题可比较凑整、拆分或商不变等方法。}
\MethodTitle{第1小题：$2.3+0.6$}

\MethodBlock{口算思路一：按相同计数单位计算}{先看成若干个十分之一，按整数加减，最后还原成小数。}{\begin{aligned}[t]&\phantom{={}}2.3+0.6\\
&=(23+6)\div 10\\
&=29\div 10\\
&=2.9\end{aligned}}
\MethodTitle{第2小题：$5.3+3.2$}

\MethodBlock{口算思路一：按相同计数单位计算}{先看成若干个十分之一，按整数加减，最后还原成小数。}{\begin{aligned}[t]&\phantom{={}}5.3+3.2\\
&=(53+32)\div 10\\
&=85\div 10\\
&=8.5\end{aligned}}
\MethodTitle{第3小题：$7.5-3.5$}

\MethodBlock{口算思路一：按相同计数单位计算}{先看成若干个十分之一，按整数加减，最后还原成小数。}{\begin{aligned}[t]&\phantom{={}}7.5-3.5\\
&=(75-35)\div 10\\
&=40\div 10\\
&=4\end{aligned}}

\MethodBlock{口算思路二：补差法}{从减数先补到一个整数，再补到被减数，两段差相加。}{\begin{aligned}[t]&\phantom{={}}7.5-3.5\\
&=(4-3.5)+(7.5-4)\\
&=0.5+3.5\\
&=4\end{aligned}}
\MethodTitle{第4小题：$1.4-0.5$}

\MethodBlock{口算思路一：按相同计数单位计算}{先看成若干个十分之一，按整数加减，最后还原成小数。}{\begin{aligned}[t]&\phantom{={}}1.4-0.5\\
&=(14-5)\div 10\\
&=9\div 10\\
&=0.9\end{aligned}}

\MethodBlock{口算思路二：补差法}{从减数先补到一个整数，再补到被减数，两段差相加。}{\begin{aligned}[t]&\phantom{={}}1.4-0.5\\
&=(1-0.5)+(1.4-1)\\
&=0.5+0.4\\
&=0.9\end{aligned}}
\MethodTitle{第5小题：$1-0.2$}

\MethodBlock{口算思路一：按相同计数单位计算}{先看成若干个十分之一，按整数加减，最后还原成小数。}{\begin{aligned}[t]&\phantom{={}}1-0.2\\
&=(10-2)\div 10\\
&=8\div 10\\
&=0.8\end{aligned}}
\MethodTitle{第6小题：$3.5-3.1$}

\MethodBlock{口算思路一：按相同计数单位计算}{先看成若干个十分之一，按整数加减，最后还原成小数。}{\begin{aligned}[t]&\phantom{={}}3.5-3.1\\
&=(35-31)\div 10\\
&=4\div 10\\
&=0.4\end{aligned}}
\MethodTitle{第7小题：$0.54+1.44$}

\MethodBlock{口算思路一：按相同计数单位计算}{先看成若干个百分之一，按整数加减，最后还原成小数。}{\begin{aligned}[t]&\phantom{={}}0.54+1.44\\
&=(54+144)\div 100\\
&=198\div 100\\
&=1.98\end{aligned}}
\MethodTitle{第8小题：$0.2+1.8$}

\MethodBlock{口算思路一：按相同计数单位计算}{先看成若干个十分之一，按整数加减，最后还原成小数。}{\begin{aligned}[t]&\phantom{={}}0.2+1.8\\
&=(2+18)\div 10\\
&=20\div 10\\
&=2\end{aligned}}

\MethodBlock{口算思路二：凑整再减去多加的数}{把第二个数凑成整数，再减去多加的部分。}{\begin{aligned}[t]&\phantom{={}}0.2+1.8\\
&=0.2+2-0.2\\
&=2.2-0.2\\
&=2\end{aligned}}
\MethodTitle{第9小题：$8.8-8.7$}

\MethodBlock{口算思路一：按相同计数单位计算}{先看成若干个十分之一，按整数加减，最后还原成小数。}{\begin{aligned}[t]&\phantom{={}}8.8-8.7\\
&=(88-87)\div 10\\
&=1\div 10\\
&=0.1\end{aligned}}
\MethodTitle{第10小题：$3.26+5.41$}

\MethodBlock{口算思路一：按相同计数单位计算}{先看成若干个百分之一，按整数加减，最后还原成小数。}{\begin{aligned}[t]&\phantom{={}}3.26+5.41\\
&=(326+541)\div 100\\
&=867\div 100\\
&=8.67\end{aligned}}
\MethodTitle{第11小题：$7.58+2.35$}

\MethodBlock{口算思路一：按相同计数单位计算}{先看成若干个百分之一，按整数加减，最后还原成小数。}{\begin{aligned}[t]&\phantom{={}}7.58+2.35\\
&=(758+235)\div 100\\
&=993\div 100\\
&=9.93\end{aligned}}
\MethodTitle{第12小题：$9.87-4.52$}

\MethodBlock{口算思路一：按相同计数单位计算}{先看成若干个百分之一，按整数加减，最后还原成小数。}{\begin{aligned}[t]&\phantom{={}}9.87-4.52\\
&=(987-452)\div 100\\
&=535\div 100\\
&=5.35\end{aligned}}

\MethodBlock{口算思路二：补差法}{从减数先补到一个整数，再补到被减数，两段差相加。}{\begin{aligned}[t]&\phantom{={}}9.87-4.52\\
&=(5-4.52)+(9.87-5)\\
&=0.48+4.87\\
&=5.35\end{aligned}}


\par\noindent

\QHeading{第2题　列竖式计算，带※的要验算。}

\par\noindent\makebox[\linewidth][c]{
\begin{tikzpicture}[x=1bp,y=-1bp,color=black,draw=black,text=black]
\begin{scope}\clip (64.0,243.0) rectangle (534.0,570.0);
\hdhead{61}{小数加减法——信息窗1}
\hdsection{1}{140}{口算。}
\hdsection{2}{325}{列竖式计算，带※的要验算。}
\hdsection{3}{947}{解决问题。}
\q{104}{185}{180}{12.5}{$2.3+0.6=$}
\q{104}{227}{180}{12.5}{$1-0.2=$}
\q{104}{270}{180}{12.5}{$8.8-8.7=$}
\q{297}{185}{180}{12.5}{$5.3+3.2=$}
\q{297}{227}{180}{12.5}{$3.5-3.1=$}
\q{297}{270}{180}{12.5}{$3.26+5.41=$}
\q{489}{185}{180}{12.5}{$7.5-3.5=$}
\q{489}{227}{180}{12.5}{$0.54+1.44=$}
\q{489}{270}{180}{12.5}{$7.58+2.35=$}
\q{676}{185}{180}{12.5}{$1.4-0.5=$}
\q{676}{227}{180}{12.5}{$0.2+1.8=$}
\q{676}{270}{180}{12.5}{$9.87-4.52=$}
\q{104}{373}{235}{12.5}{$3.62+5.78=$}
\q{364}{373}{235}{12.5}{$2.3-0.84=$}
\q{615}{373}{235}{12.5}{※$6.3+2.96=$}
\q{104}{561}{235}{12.5}{$5-3.8=$}
\q{364}{561}{235}{12.5}{$4.2+3.18=$}
\q{615}{561}{235}{12.5}{※$9.4-7.56=$}
\q{104}{748}{235}{12.5}{$47.1-4.71=$}
\q{364}{748}{235}{12.5}{$5.79+23.6=$}
\q{615}{748}{235}{12.5}{※$30-9.29=$}
\q{105}{995}{715}{12.5}{小刚要买一个16.85元的文具盒，他带了20元，买完文具盒后还剩多}
\q{105}{1038}{715}{12.5}{少钱？}

\begin{scope}[color=answerred,draw=answerred,text=answerred]\ans{143.33}{121.33}{46.67}{10.5}{$2.9$}
\ans{143.33}{149.33}{46.67}{10.5}{$0.8$}
\ans{143.33}{177.33}{46.67}{10.5}{$0.1$}
\ans{288.00}{121.33}{46.67}{10.5}{$8.5$}
\ans{288.00}{149.33}{46.67}{10.5}{$0.4$}
\ans{288.00}{177.33}{46.67}{10.5}{$8.67$}
\ans{410.67}{121.33}{46.67}{10.5}{$4$}
\ans{410.67}{149.33}{46.67}{10.5}{$1.98$}
\ans{410.67}{177.33}{46.67}{10.5}{$9.93$}
\ans{536.67}{121.33}{46.67}{10.5}{$0.9$}
\ans{536.67}{149.33}{46.67}{10.5}{$2$}
\ans{536.67}{177.33}{46.67}{10.5}{$5.35$}
\ans{69.33}{276.00}{142.67}{11}{\begin{tikzpicture}[x=1em,y=1em,baseline=(current bounding box.center),draw=answerred,text=answerred]
\node[anchor=base,inner sep=0pt] at (0.000,0) {$3$};
\node[anchor=base,inner sep=0pt] at (0.800,0) {$.$};
\node[anchor=base,inner sep=0pt] at (1.100,0) {$6$};
\node[anchor=base,inner sep=0pt] at (1.900,0) {$2$};
\node[anchor=base,inner sep=0pt] at (0.000,-1.1) {$5$};
\node[anchor=base,inner sep=0pt] at (0.800,-1.1) {$.$};
\node[anchor=base,inner sep=0pt] at (1.100,-1.1) {$7$};
\node[anchor=base,inner sep=0pt] at (1.900,-1.1) {$8$};
\node[anchor=base,inner sep=0pt] at (0.000,-2.3) {$9$};
\node[anchor=base,inner sep=0pt] at (0.800,-2.3) {$.$};
\node[anchor=base,inner sep=0pt] at (1.100,-2.3) {$4$};
\node[anchor=base,inner sep=0pt] at (1.900,-2.3) {$0$};
\node[anchor=base,inner sep=0pt] at (-.7,-1.1) {$+$};
\draw[line width=.5pt] (-1,-1.52)--(2.500,-1.52);
\node[anchor=base,inner sep=0pt,font=\normalsize] at (1.100,.72) {$1$};
\node[anchor=base,inner sep=0pt,font=\normalsize] at (0.000,.72) {$1$};
\end{tikzpicture}}
\ans{242.67}{276.00}{142.67}{11}{\begin{tikzpicture}[x=1em,y=1em,baseline=(current bounding box.center),draw=answerred,text=answerred]
\node[anchor=base,inner sep=0pt] at (0.000,0) {$2$};
\node[anchor=base,inner sep=0pt] at (0.800,0) {$.$};
\node[anchor=base,inner sep=0pt] at (1.100,0) {$3$};
\node[anchor=base,inner sep=0pt] at (1.900,0) {$0$};
\node[anchor=base,inner sep=0pt] at (0.000,-1.1) {$0$};
\node[anchor=base,inner sep=0pt] at (0.800,-1.1) {$.$};
\node[anchor=base,inner sep=0pt] at (1.100,-1.1) {$8$};
\node[anchor=base,inner sep=0pt] at (1.900,-1.1) {$4$};
\node[anchor=base,inner sep=0pt] at (0.000,-2.3) {$1$};
\node[anchor=base,inner sep=0pt] at (0.800,-2.3) {$.$};
\node[anchor=base,inner sep=0pt] at (1.100,-2.3) {$4$};
\node[anchor=base,inner sep=0pt] at (1.900,-2.3) {$6$};
\node[anchor=base,inner sep=0pt] at (-.7,-1.1) {$-$};
\draw[line width=.5pt] (-1,-1.52)--(2.500,-1.52);
\draw[line width=.35pt] (-0.160,.03)--(0.160,.43);
\node[anchor=base,inner sep=0pt,font=\normalsize] at (0.000,.72) {$1$};
\draw[line width=.35pt] (0.940,.03)--(1.260,.43);
\node[anchor=base,inner sep=0pt,font=\normalsize] at (1.100,.72) {$12$};
\draw[line width=.35pt] (1.740,.03)--(2.060,.43);
\node[anchor=base,inner sep=0pt,font=\normalsize] at (1.900,.72) {$10$};
\end{tikzpicture}}
\ans{410.00}{276.00}{142.67}{9.8}{\begin{tikzpicture}[x=1em,y=1em,baseline=(current bounding box.center),draw=answerred,text=answerred]
\node[anchor=base,inner sep=0pt] at (0.000,0) {$6$};
\node[anchor=base,inner sep=0pt] at (0.800,0) {$.$};
\node[anchor=base,inner sep=0pt] at (1.100,0) {$3$};
\node[anchor=base,inner sep=0pt] at (1.900,0) {$0$};
\node[anchor=base,inner sep=0pt] at (0.000,-1.1) {$2$};
\node[anchor=base,inner sep=0pt] at (0.800,-1.1) {$.$};
\node[anchor=base,inner sep=0pt] at (1.100,-1.1) {$9$};
\node[anchor=base,inner sep=0pt] at (1.900,-1.1) {$6$};
\node[anchor=base,inner sep=0pt] at (0.000,-2.3) {$9$};
\node[anchor=base,inner sep=0pt] at (0.800,-2.3) {$.$};
\node[anchor=base,inner sep=0pt] at (1.100,-2.3) {$2$};
\node[anchor=base,inner sep=0pt] at (1.900,-2.3) {$6$};
\node[anchor=base,inner sep=0pt] at (-.7,-1.1) {$+$};
\draw[line width=.5pt] (-1,-1.52)--(2.500,-1.52);
\node[anchor=base,inner sep=0pt,font=\normalsize] at (0.000,.72) {$1$};
\end{tikzpicture}\hspace{3pt}验算：\begin{tikzpicture}[x=1em,y=1em,baseline=(current bounding box.center),draw=answerred,text=answerred]
\node[anchor=base,inner sep=0pt] at (0.000,0) {$9$};
\node[anchor=base,inner sep=0pt] at (0.800,0) {$.$};
\node[anchor=base,inner sep=0pt] at (1.100,0) {$2$};
\node[anchor=base,inner sep=0pt] at (1.900,0) {$6$};
\node[anchor=base,inner sep=0pt] at (0.000,-1.1) {$2$};
\node[anchor=base,inner sep=0pt] at (0.800,-1.1) {$.$};
\node[anchor=base,inner sep=0pt] at (1.100,-1.1) {$9$};
\node[anchor=base,inner sep=0pt] at (1.900,-1.1) {$6$};
\node[anchor=base,inner sep=0pt] at (0.000,-2.3) {$6$};
\node[anchor=base,inner sep=0pt] at (0.800,-2.3) {$.$};
\node[anchor=base,inner sep=0pt] at (1.100,-2.3) {$3$};
\node[anchor=base,inner sep=0pt] at (1.900,-2.3) {$0$};
\node[anchor=base,inner sep=0pt] at (-.7,-1.1) {$-$};
\draw[line width=.5pt] (-1,-1.52)--(2.500,-1.52);
\draw[line width=.35pt] (-0.160,.03)--(0.160,.43);
\node[anchor=base,inner sep=0pt,font=\normalsize] at (0.000,.72) {$8$};
\draw[line width=.35pt] (0.940,.03)--(1.260,.43);
\node[anchor=base,inner sep=0pt,font=\normalsize] at (1.100,.72) {$12$};
\end{tikzpicture}}
\ans{69.33}{402.67}{142.67}{11}{\begin{tikzpicture}[x=1em,y=1em,baseline=(current bounding box.center),draw=answerred,text=answerred]
\node[anchor=base,inner sep=0pt] at (0.000,0) {$5$};
\node[anchor=base,inner sep=0pt] at (0.800,0) {$.$};
\node[anchor=base,inner sep=0pt] at (1.100,0) {$0$};
\node[anchor=base,inner sep=0pt] at (0.000,-1.1) {$3$};
\node[anchor=base,inner sep=0pt] at (0.800,-1.1) {$.$};
\node[anchor=base,inner sep=0pt] at (1.100,-1.1) {$8$};
\node[anchor=base,inner sep=0pt] at (0.000,-2.3) {$1$};
\node[anchor=base,inner sep=0pt] at (0.800,-2.3) {$.$};
\node[anchor=base,inner sep=0pt] at (1.100,-2.3) {$2$};
\node[anchor=base,inner sep=0pt] at (-.7,-1.1) {$-$};
\draw[line width=.5pt] (-1,-1.52)--(1.700,-1.52);
\draw[line width=.35pt] (-0.160,.03)--(0.160,.43);
\node[anchor=base,inner sep=0pt,font=\normalsize] at (0.000,.72) {$4$};
\draw[line width=.35pt] (0.940,.03)--(1.260,.43);
\node[anchor=base,inner sep=0pt,font=\normalsize] at (1.100,.72) {$10$};
\end{tikzpicture}}
\ans{242.67}{402.67}{142.67}{11}{\begin{tikzpicture}[x=1em,y=1em,baseline=(current bounding box.center),draw=answerred,text=answerred]
\node[anchor=base,inner sep=0pt] at (0.000,0) {$4$};
\node[anchor=base,inner sep=0pt] at (0.800,0) {$.$};
\node[anchor=base,inner sep=0pt] at (1.100,0) {$2$};
\node[anchor=base,inner sep=0pt] at (1.900,0) {$0$};
\node[anchor=base,inner sep=0pt] at (0.000,-1.1) {$3$};
\node[anchor=base,inner sep=0pt] at (0.800,-1.1) {$.$};
\node[anchor=base,inner sep=0pt] at (1.100,-1.1) {$1$};
\node[anchor=base,inner sep=0pt] at (1.900,-1.1) {$8$};
\node[anchor=base,inner sep=0pt] at (0.000,-2.3) {$7$};
\node[anchor=base,inner sep=0pt] at (0.800,-2.3) {$.$};
\node[anchor=base,inner sep=0pt] at (1.100,-2.3) {$3$};
\node[anchor=base,inner sep=0pt] at (1.900,-2.3) {$8$};
\node[anchor=base,inner sep=0pt] at (-.7,-1.1) {$+$};
\draw[line width=.5pt] (-1,-1.52)--(2.500,-1.52);
\end{tikzpicture}}
\ans{410.00}{402.67}{142.67}{9.8}{\begin{tikzpicture}[x=1em,y=1em,baseline=(current bounding box.center),draw=answerred,text=answerred]
\node[anchor=base,inner sep=0pt] at (0.000,0) {$9$};
\node[anchor=base,inner sep=0pt] at (0.800,0) {$.$};
\node[anchor=base,inner sep=0pt] at (1.100,0) {$4$};
\node[anchor=base,inner sep=0pt] at (1.900,0) {$0$};
\node[anchor=base,inner sep=0pt] at (0.000,-1.1) {$7$};
\node[anchor=base,inner sep=0pt] at (0.800,-1.1) {$.$};
\node[anchor=base,inner sep=0pt] at (1.100,-1.1) {$5$};
\node[anchor=base,inner sep=0pt] at (1.900,-1.1) {$6$};
\node[anchor=base,inner sep=0pt] at (0.000,-2.3) {$1$};
\node[anchor=base,inner sep=0pt] at (0.800,-2.3) {$.$};
\node[anchor=base,inner sep=0pt] at (1.100,-2.3) {$8$};
\node[anchor=base,inner sep=0pt] at (1.900,-2.3) {$4$};
\node[anchor=base,inner sep=0pt] at (-.7,-1.1) {$-$};
\draw[line width=.5pt] (-1,-1.52)--(2.500,-1.52);
\draw[line width=.35pt] (-0.160,.03)--(0.160,.43);
\node[anchor=base,inner sep=0pt,font=\normalsize] at (0.000,.72) {$8$};
\draw[line width=.35pt] (0.940,.03)--(1.260,.43);
\node[anchor=base,inner sep=0pt,font=\normalsize] at (1.100,.72) {$13$};
\draw[line width=.35pt] (1.740,.03)--(2.060,.43);
\node[anchor=base,inner sep=0pt,font=\normalsize] at (1.900,.72) {$10$};
\end{tikzpicture}\hspace{3pt}验算：\begin{tikzpicture}[x=1em,y=1em,baseline=(current bounding box.center),draw=answerred,text=answerred]
\node[anchor=base,inner sep=0pt] at (0.000,0) {$1$};
\node[anchor=base,inner sep=0pt] at (0.800,0) {$.$};
\node[anchor=base,inner sep=0pt] at (1.100,0) {$8$};
\node[anchor=base,inner sep=0pt] at (1.900,0) {$4$};
\node[anchor=base,inner sep=0pt] at (0.000,-1.1) {$7$};
\node[anchor=base,inner sep=0pt] at (0.800,-1.1) {$.$};
\node[anchor=base,inner sep=0pt] at (1.100,-1.1) {$5$};
\node[anchor=base,inner sep=0pt] at (1.900,-1.1) {$6$};
\node[anchor=base,inner sep=0pt] at (0.000,-2.3) {$9$};
\node[anchor=base,inner sep=0pt] at (0.800,-2.3) {$.$};
\node[anchor=base,inner sep=0pt] at (1.100,-2.3) {$4$};
\node[anchor=base,inner sep=0pt] at (1.900,-2.3) {$0$};
\node[anchor=base,inner sep=0pt] at (-.7,-1.1) {$+$};
\draw[line width=.5pt] (-1,-1.52)--(2.500,-1.52);
\node[anchor=base,inner sep=0pt,font=\normalsize] at (1.100,.72) {$1$};
\node[anchor=base,inner sep=0pt,font=\normalsize] at (0.000,.72) {$1$};
\end{tikzpicture}}
\ans{69.33}{526.67}{142.67}{11}{\begin{tikzpicture}[x=1em,y=1em,baseline=(current bounding box.center),draw=answerred,text=answerred]
\node[anchor=base,inner sep=0pt] at (0.000,0) {$4$};
\node[anchor=base,inner sep=0pt] at (0.800,0) {$7$};
\node[anchor=base,inner sep=0pt] at (1.600,0) {$.$};
\node[anchor=base,inner sep=0pt] at (1.900,0) {$1$};
\node[anchor=base,inner sep=0pt] at (2.700,0) {$0$};
\node[anchor=base,inner sep=0pt] at (0.800,-1.1) {$4$};
\node[anchor=base,inner sep=0pt] at (1.600,-1.1) {$.$};
\node[anchor=base,inner sep=0pt] at (1.900,-1.1) {$7$};
\node[anchor=base,inner sep=0pt] at (2.700,-1.1) {$1$};
\node[anchor=base,inner sep=0pt] at (0.000,-2.3) {$4$};
\node[anchor=base,inner sep=0pt] at (0.800,-2.3) {$2$};
\node[anchor=base,inner sep=0pt] at (1.600,-2.3) {$.$};
\node[anchor=base,inner sep=0pt] at (1.900,-2.3) {$3$};
\node[anchor=base,inner sep=0pt] at (2.700,-2.3) {$9$};
\node[anchor=base,inner sep=0pt] at (-.7,-1.1) {$-$};
\draw[line width=.5pt] (-1,-1.52)--(3.300,-1.52);
\draw[line width=.35pt] (0.640,.03)--(0.960,.43);
\node[anchor=base,inner sep=0pt,font=\normalsize] at (0.800,.72) {$6$};
\draw[line width=.35pt] (1.740,.03)--(2.060,.43);
\node[anchor=base,inner sep=0pt,font=\normalsize] at (1.900,.72) {$10$};
\draw[line width=.35pt] (2.540,.03)--(2.860,.43);
\node[anchor=base,inner sep=0pt,font=\normalsize] at (2.700,.72) {$10$};
\end{tikzpicture}}
\ans{242.67}{526.67}{142.67}{11}{\begin{tikzpicture}[x=1em,y=1em,baseline=(current bounding box.center),draw=answerred,text=answerred]
\node[anchor=base,inner sep=0pt] at (0.800,0) {$5$};
\node[anchor=base,inner sep=0pt] at (1.600,0) {$.$};
\node[anchor=base,inner sep=0pt] at (1.900,0) {$7$};
\node[anchor=base,inner sep=0pt] at (2.700,0) {$9$};
\node[anchor=base,inner sep=0pt] at (0.000,-1.1) {$2$};
\node[anchor=base,inner sep=0pt] at (0.800,-1.1) {$3$};
\node[anchor=base,inner sep=0pt] at (1.600,-1.1) {$.$};
\node[anchor=base,inner sep=0pt] at (1.900,-1.1) {$6$};
\node[anchor=base,inner sep=0pt] at (2.700,-1.1) {$0$};
\node[anchor=base,inner sep=0pt] at (0.000,-2.3) {$2$};
\node[anchor=base,inner sep=0pt] at (0.800,-2.3) {$9$};
\node[anchor=base,inner sep=0pt] at (1.600,-2.3) {$.$};
\node[anchor=base,inner sep=0pt] at (1.900,-2.3) {$3$};
\node[anchor=base,inner sep=0pt] at (2.700,-2.3) {$9$};
\node[anchor=base,inner sep=0pt] at (-.7,-1.1) {$+$};
\draw[line width=.5pt] (-1,-1.52)--(3.300,-1.52);
\node[anchor=base,inner sep=0pt,font=\normalsize] at (0.800,.72) {$1$};
\end{tikzpicture}}
\ans{410.00}{526.67}{142.67}{9.8}{\begin{tikzpicture}[x=1em,y=1em,baseline=(current bounding box.center),draw=answerred,text=answerred]
\node[anchor=base,inner sep=0pt] at (0.000,0) {$3$};
\node[anchor=base,inner sep=0pt] at (0.800,0) {$0$};
\node[anchor=base,inner sep=0pt] at (1.600,0) {$.$};
\node[anchor=base,inner sep=0pt] at (1.900,0) {$0$};
\node[anchor=base,inner sep=0pt] at (2.700,0) {$0$};
\node[anchor=base,inner sep=0pt] at (0.800,-1.1) {$9$};
\node[anchor=base,inner sep=0pt] at (1.600,-1.1) {$.$};
\node[anchor=base,inner sep=0pt] at (1.900,-1.1) {$2$};
\node[anchor=base,inner sep=0pt] at (2.700,-1.1) {$9$};
\node[anchor=base,inner sep=0pt] at (0.000,-2.3) {$2$};
\node[anchor=base,inner sep=0pt] at (0.800,-2.3) {$0$};
\node[anchor=base,inner sep=0pt] at (1.600,-2.3) {$.$};
\node[anchor=base,inner sep=0pt] at (1.900,-2.3) {$7$};
\node[anchor=base,inner sep=0pt] at (2.700,-2.3) {$1$};
\node[anchor=base,inner sep=0pt] at (-.7,-1.1) {$-$};
\draw[line width=.5pt] (-1,-1.52)--(3.300,-1.52);
\draw[line width=.35pt] (-0.160,.03)--(0.160,.43);
\node[anchor=base,inner sep=0pt,font=\normalsize] at (0.000,.72) {$2$};
\draw[line width=.35pt] (0.640,.03)--(0.960,.43);
\node[anchor=base,inner sep=0pt,font=\normalsize] at (0.800,.72) {$9$};
\draw[line width=.35pt] (1.740,.03)--(2.060,.43);
\node[anchor=base,inner sep=0pt,font=\normalsize] at (1.900,.72) {$9$};
\draw[line width=.35pt] (2.540,.03)--(2.860,.43);
\node[anchor=base,inner sep=0pt,font=\normalsize] at (2.700,.72) {$10$};
\end{tikzpicture}\hspace{3pt}验算：\begin{tikzpicture}[x=1em,y=1em,baseline=(current bounding box.center),draw=answerred,text=answerred]
\node[anchor=base,inner sep=0pt] at (0.000,0) {$2$};
\node[anchor=base,inner sep=0pt] at (0.800,0) {$0$};
\node[anchor=base,inner sep=0pt] at (1.600,0) {$.$};
\node[anchor=base,inner sep=0pt] at (1.900,0) {$7$};
\node[anchor=base,inner sep=0pt] at (2.700,0) {$1$};
\node[anchor=base,inner sep=0pt] at (0.800,-1.1) {$9$};
\node[anchor=base,inner sep=0pt] at (1.600,-1.1) {$.$};
\node[anchor=base,inner sep=0pt] at (1.900,-1.1) {$2$};
\node[anchor=base,inner sep=0pt] at (2.700,-1.1) {$9$};
\node[anchor=base,inner sep=0pt] at (0.000,-2.3) {$3$};
\node[anchor=base,inner sep=0pt] at (0.800,-2.3) {$0$};
\node[anchor=base,inner sep=0pt] at (1.600,-2.3) {$.$};
\node[anchor=base,inner sep=0pt] at (1.900,-2.3) {$0$};
\node[anchor=base,inner sep=0pt] at (2.700,-2.3) {$0$};
\node[anchor=base,inner sep=0pt] at (-.7,-1.1) {$+$};
\draw[line width=.5pt] (-1,-1.52)--(3.300,-1.52);
\node[anchor=base,inner sep=0pt,font=\normalsize] at (1.900,.72) {$1$};
\node[anchor=base,inner sep=0pt,font=\normalsize] at (0.800,.72) {$1$};
\node[anchor=base,inner sep=0pt,font=\normalsize] at (0.000,.72) {$1$};
\end{tikzpicture}}
\ans{72.67}{720.67}{454.00}{11}{$20-16.85=3.15\text{（元）}$\\答：还剩3.15元。}


% 本次补齐原书第61页第2题横式等号后的得数。
\ans{144.1199}{246.7304}{43}{12.5}{$\vphantom{0}9.4$}
\ans{311.5588}{246.7304}{43}{12.5}{$\vphantom{0}1.46$}
\ans{488.1799}{246.7304}{43}{12.5}{$\vphantom{0}9.26$}
\ans{122.6512}{372.0634}{43}{12.5}{$\vphantom{0}1.2$}
\ans{311.3568}{372.0634}{43}{12.5}{$\vphantom{0}7.38$}
\ans{488.3819}{372.0634}{43}{12.5}{$\vphantom{0}1.84$}
\ans{144.3219}{496.7304}{43}{12.5}{$\vphantom{0}42.39$}
\ans{317.4602}{496.7304}{43}{12.5}{$\vphantom{0}29.39$}
\ans{484.9948}{496.7304}{43}{12.5}{$\vphantom{0}20.71$}
\end{scope}
\end{scope}\pgfresetboundingbox\path[use as bounding box] (64.0,243.0) rectangle (534.0,570.0);
\end{tikzpicture}}\par



\par






\par\noindent

\Answer{
\QHeading{计算过程：2. 列竖式计算}

小数点对齐，缺位补零；从右往左逐列计算。下面完整列出进位、退位及各数位的计算过程。

\par\medskip\noindent
\textbf{第1小题：\(3.62+5.78=9.4\)}\par\smallskip
\noindent
\begin{tikzpicture}[x=1em,y=1em,baseline=(current bounding box.center),draw=answerred,text=answerred]
\node[anchor=base,inner sep=0pt] at (0.000,0) {$3$};
\node[anchor=base,inner sep=0pt] at (0.800,0) {$.$};
\node[anchor=base,inner sep=0pt] at (1.100,0) {$6$};
\node[anchor=base,inner sep=0pt] at (1.900,0) {$2$};
\node[anchor=base,inner sep=0pt] at (0.000,-1.1) {$5$};
\node[anchor=base,inner sep=0pt] at (0.800,-1.1) {$.$};
\node[anchor=base,inner sep=0pt] at (1.100,-1.1) {$7$};
\node[anchor=base,inner sep=0pt] at (1.900,-1.1) {$8$};
\node[anchor=base,inner sep=0pt] at (0.000,-2.3) {$9$};
\node[anchor=base,inner sep=0pt] at (0.800,-2.3) {$.$};
\node[anchor=base,inner sep=0pt] at (1.100,-2.3) {$4$};
\node[anchor=base,inner sep=0pt] at (1.900,-2.3) {$0$};
\node[anchor=base,inner sep=0pt] at (-.7,-1.1) {$+$};
\draw[line width=.5pt] (-1,-1.52)--(2.500,-1.52);
\node[anchor=base,inner sep=0pt,font=\normalsize] at (1.100,.72) {$1$};
\node[anchor=base,inner sep=0pt,font=\normalsize] at (0.000,.72) {$1$};
\end{tikzpicture}
\par\smallskip{\normalsize 上方小字表示从后一位进来的数。\par}
\vspace{0pt}
\begin{enumerate}[leftmargin=1.7em,itemsep=5pt,topsep=0pt]
\item 百分位：\(\begin{aligned}[t]&\phantom{={}}2+8\\&=10\end{aligned}\)\quad 本位写0，向前一位进1。
\item 十分位：\(\begin{aligned}[t]&\phantom{={}}6+7+1\\&=13+1\\&=14\end{aligned}\)\quad 本位写4，向前一位进1。
\item 个位：\(\begin{aligned}[t]&\phantom{={}}3+5+1\\&=8+1\\&=9\end{aligned}\)\quad 本位写9。
\end{enumerate}
\par\smallskip
因此，\(\begin{aligned}[t]&\phantom{={}}3.62+5.78\\&=9.40\end{aligned}\)。
\par\medskip

\par\medskip\noindent
\textbf{第2小题：\(2.3-0.84=1.46\)}\par\smallskip
\noindent
\begin{tikzpicture}[x=1em,y=1em,baseline=(current bounding box.center),draw=answerred,text=answerred]
\node[anchor=base,inner sep=0pt] at (0.000,0) {$2$};
\node[anchor=base,inner sep=0pt] at (0.800,0) {$.$};
\node[anchor=base,inner sep=0pt] at (1.100,0) {$3$};
\node[anchor=base,inner sep=0pt] at (1.900,0) {$0$};
\node[anchor=base,inner sep=0pt] at (0.000,-1.1) {$0$};
\node[anchor=base,inner sep=0pt] at (0.800,-1.1) {$.$};
\node[anchor=base,inner sep=0pt] at (1.100,-1.1) {$8$};
\node[anchor=base,inner sep=0pt] at (1.900,-1.1) {$4$};
\node[anchor=base,inner sep=0pt] at (0.000,-2.3) {$1$};
\node[anchor=base,inner sep=0pt] at (0.800,-2.3) {$.$};
\node[anchor=base,inner sep=0pt] at (1.100,-2.3) {$4$};
\node[anchor=base,inner sep=0pt] at (1.900,-2.3) {$6$};
\node[anchor=base,inner sep=0pt] at (-.7,-1.1) {$-$};
\draw[line width=.5pt] (-1,-1.52)--(2.500,-1.52);
\draw[line width=.35pt] (-0.160,.03)--(0.160,.43);
\node[anchor=base,inner sep=0pt,font=\normalsize] at (0.000,.72) {$1$};
\draw[line width=.35pt] (0.940,.03)--(1.260,.43);
\node[anchor=base,inner sep=0pt,font=\normalsize] at (1.100,.72) {$12$};
\draw[line width=.35pt] (1.740,.03)--(2.060,.43);
\node[anchor=base,inner sep=0pt,font=\normalsize] at (1.900,.72) {$10$};
\end{tikzpicture}
\par\smallskip{\normalsize 划线数字是退位前的数，上方小字是退位后参与计算的数。\par}
\vspace{0pt}
\begin{enumerate}[leftmargin=1.7em,itemsep=5pt,topsep=0pt]
\item 十分位退位：\(\begin{aligned}[t]&\phantom{={}}3-1\\&=2\end{aligned}\)\quad 借出1给后一位。
\item 百分位接收退位：\(\begin{aligned}[t]&\phantom{={}}0+10\\&=10\end{aligned}\)\quad 前一位的1变成本位的10。
\item 百分位：\(\begin{aligned}[t]&\phantom{={}}10-4\\&=6\end{aligned}\)\quad 本位写6。
\item 个位退位：\(\begin{aligned}[t]&\phantom{={}}2-1\\&=1\end{aligned}\)\quad 借出1给后一位。
\item 十分位接收退位：\(\begin{aligned}[t]&\phantom{={}}2+10\\&=12\end{aligned}\)\quad 前一位的1变成本位的10。
\item 十分位：\(\begin{aligned}[t]&\phantom{={}}12-8\\&=4\end{aligned}\)\quad 本位写4。
\item 个位：\(\begin{aligned}[t]&\phantom{={}}1-0\\&=1\end{aligned}\)\quad 本位写1。
\end{enumerate}
\par\smallskip
因此，\(\begin{aligned}[t]&\phantom{={}}2.3-0.84\\&=1.46\end{aligned}\)。
\par\medskip

\par\medskip\noindent
\textbf{第3小题：\(6.3+2.96=9.26\)}\par\smallskip
\noindent
\begin{tikzpicture}[x=1em,y=1em,baseline=(current bounding box.center),draw=answerred,text=answerred]
\node[anchor=base,inner sep=0pt] at (0.000,0) {$6$};
\node[anchor=base,inner sep=0pt] at (0.800,0) {$.$};
\node[anchor=base,inner sep=0pt] at (1.100,0) {$3$};
\node[anchor=base,inner sep=0pt] at (1.900,0) {$0$};
\node[anchor=base,inner sep=0pt] at (0.000,-1.1) {$2$};
\node[anchor=base,inner sep=0pt] at (0.800,-1.1) {$.$};
\node[anchor=base,inner sep=0pt] at (1.100,-1.1) {$9$};
\node[anchor=base,inner sep=0pt] at (1.900,-1.1) {$6$};
\node[anchor=base,inner sep=0pt] at (0.000,-2.3) {$9$};
\node[anchor=base,inner sep=0pt] at (0.800,-2.3) {$.$};
\node[anchor=base,inner sep=0pt] at (1.100,-2.3) {$2$};
\node[anchor=base,inner sep=0pt] at (1.900,-2.3) {$6$};
\node[anchor=base,inner sep=0pt] at (-.7,-1.1) {$+$};
\draw[line width=.5pt] (-1,-1.52)--(2.500,-1.52);
\node[anchor=base,inner sep=0pt,font=\normalsize] at (0.000,.72) {$1$};
\end{tikzpicture}
\par\smallskip{\normalsize 上方小字表示从后一位进来的数。\par}
\vspace{0pt}
\begin{enumerate}[leftmargin=1.7em,itemsep=5pt,topsep=0pt]
\item 百分位：\(\begin{aligned}[t]&\phantom{={}}0+6\\&=6\end{aligned}\)\quad 本位写6。
\item 十分位：\(\begin{aligned}[t]&\phantom{={}}3+9\\&=12\end{aligned}\)\quad 本位写2，向前一位进1。
\item 个位：\(\begin{aligned}[t]&\phantom{={}}6+2+1\\&=8+1\\&=9\end{aligned}\)\quad 本位写9。
\end{enumerate}
\par\smallskip
因此，\(\begin{aligned}[t]&\phantom{={}}6.3+2.96\\&=9.26\end{aligned}\)。
\par\medskip

\par\medskip\noindent
\textbf{第3小题验算：\(9.26-2.96\)}\par\smallskip
\noindent
\begin{tikzpicture}[x=1em,y=1em,baseline=(current bounding box.center),draw=answerred,text=answerred]
\node[anchor=base,inner sep=0pt] at (0.000,0) {$9$};
\node[anchor=base,inner sep=0pt] at (0.800,0) {$.$};
\node[anchor=base,inner sep=0pt] at (1.100,0) {$2$};
\node[anchor=base,inner sep=0pt] at (1.900,0) {$6$};
\node[anchor=base,inner sep=0pt] at (0.000,-1.1) {$2$};
\node[anchor=base,inner sep=0pt] at (0.800,-1.1) {$.$};
\node[anchor=base,inner sep=0pt] at (1.100,-1.1) {$9$};
\node[anchor=base,inner sep=0pt] at (1.900,-1.1) {$6$};
\node[anchor=base,inner sep=0pt] at (0.000,-2.3) {$6$};
\node[anchor=base,inner sep=0pt] at (0.800,-2.3) {$.$};
\node[anchor=base,inner sep=0pt] at (1.100,-2.3) {$3$};
\node[anchor=base,inner sep=0pt] at (1.900,-2.3) {$0$};
\node[anchor=base,inner sep=0pt] at (-.7,-1.1) {$-$};
\draw[line width=.5pt] (-1,-1.52)--(2.500,-1.52);
\draw[line width=.35pt] (-0.160,.03)--(0.160,.43);
\node[anchor=base,inner sep=0pt,font=\normalsize] at (0.000,.72) {$8$};
\draw[line width=.35pt] (0.940,.03)--(1.260,.43);
\node[anchor=base,inner sep=0pt,font=\normalsize] at (1.100,.72) {$12$};
\end{tikzpicture}
\par\smallskip{\normalsize 划线数字是退位前的数，上方小字是退位后参与计算的数。\par}
\vspace{0pt}
\begin{enumerate}[leftmargin=1.7em,itemsep=5pt,topsep=0pt]
\item 百分位：\(\begin{aligned}[t]&\phantom{={}}6-6\\&=0\end{aligned}\)\quad 本位写0。
\item 个位退位：\(\begin{aligned}[t]&\phantom{={}}9-1\\&=8\end{aligned}\)\quad 借出1给后一位。
\item 十分位接收退位：\(\begin{aligned}[t]&\phantom{={}}2+10\\&=12\end{aligned}\)\quad 前一位的1变成本位的10。
\item 十分位：\(\begin{aligned}[t]&\phantom{={}}12-9\\&=3\end{aligned}\)\quad 本位写3。
\item 个位：\(\begin{aligned}[t]&\phantom{={}}8-2\\&=6\end{aligned}\)\quad 本位写6。
\end{enumerate}
\par\smallskip
因此，\(\begin{aligned}[t]&\phantom{={}}9.26-2.96\\&=6.30\end{aligned}\)。
\par\medskip

\par\medskip\noindent
\textbf{第4小题：\(5-3.8=1.2\)}\par\smallskip
\noindent
\begin{tikzpicture}[x=1em,y=1em,baseline=(current bounding box.center),draw=answerred,text=answerred]
\node[anchor=base,inner sep=0pt] at (0.000,0) {$5$};
\node[anchor=base,inner sep=0pt] at (0.800,0) {$.$};
\node[anchor=base,inner sep=0pt] at (1.100,0) {$0$};
\node[anchor=base,inner sep=0pt] at (0.000,-1.1) {$3$};
\node[anchor=base,inner sep=0pt] at (0.800,-1.1) {$.$};
\node[anchor=base,inner sep=0pt] at (1.100,-1.1) {$8$};
\node[anchor=base,inner sep=0pt] at (0.000,-2.3) {$1$};
\node[anchor=base,inner sep=0pt] at (0.800,-2.3) {$.$};
\node[anchor=base,inner sep=0pt] at (1.100,-2.3) {$2$};
\node[anchor=base,inner sep=0pt] at (-.7,-1.1) {$-$};
\draw[line width=.5pt] (-1,-1.52)--(1.700,-1.52);
\draw[line width=.35pt] (-0.160,.03)--(0.160,.43);
\node[anchor=base,inner sep=0pt,font=\normalsize] at (0.000,.72) {$4$};
\draw[line width=.35pt] (0.940,.03)--(1.260,.43);
\node[anchor=base,inner sep=0pt,font=\normalsize] at (1.100,.72) {$10$};
\end{tikzpicture}
\par\smallskip{\normalsize 划线数字是退位前的数，上方小字是退位后参与计算的数。\par}
\vspace{0pt}
\begin{enumerate}[leftmargin=1.7em,itemsep=5pt,topsep=0pt]
\item 个位退位：\(\begin{aligned}[t]&\phantom{={}}5-1\\&=4\end{aligned}\)\quad 借出1给后一位。
\item 十分位接收退位：\(\begin{aligned}[t]&\phantom{={}}0+10\\&=10\end{aligned}\)\quad 前一位的1变成本位的10。
\item 十分位：\(\begin{aligned}[t]&\phantom{={}}10-8\\&=2\end{aligned}\)\quad 本位写2。
\item 个位：\(\begin{aligned}[t]&\phantom{={}}4-3\\&=1\end{aligned}\)\quad 本位写1。
\end{enumerate}
\par\smallskip
因此，\(\begin{aligned}[t]&\phantom{={}}5-3.8\\&=1.2\end{aligned}\)。
\par\medskip

\par\medskip\noindent
\textbf{第5小题：\(4.2+3.18=7.38\)}\par\smallskip
\noindent
\begin{tikzpicture}[x=1em,y=1em,baseline=(current bounding box.center),draw=answerred,text=answerred]
\node[anchor=base,inner sep=0pt] at (0.000,0) {$4$};
\node[anchor=base,inner sep=0pt] at (0.800,0) {$.$};
\node[anchor=base,inner sep=0pt] at (1.100,0) {$2$};
\node[anchor=base,inner sep=0pt] at (1.900,0) {$0$};
\node[anchor=base,inner sep=0pt] at (0.000,-1.1) {$3$};
\node[anchor=base,inner sep=0pt] at (0.800,-1.1) {$.$};
\node[anchor=base,inner sep=0pt] at (1.100,-1.1) {$1$};
\node[anchor=base,inner sep=0pt] at (1.900,-1.1) {$8$};
\node[anchor=base,inner sep=0pt] at (0.000,-2.3) {$7$};
\node[anchor=base,inner sep=0pt] at (0.800,-2.3) {$.$};
\node[anchor=base,inner sep=0pt] at (1.100,-2.3) {$3$};
\node[anchor=base,inner sep=0pt] at (1.900,-2.3) {$8$};
\node[anchor=base,inner sep=0pt] at (-.7,-1.1) {$+$};
\draw[line width=.5pt] (-1,-1.52)--(2.500,-1.52);
\end{tikzpicture}
\par\smallskip{\normalsize 上方小字表示从后一位进来的数。\par}
\vspace{0pt}
\begin{enumerate}[leftmargin=1.7em,itemsep=5pt,topsep=0pt]
\item 百分位：\(\begin{aligned}[t]&\phantom{={}}0+8\\&=8\end{aligned}\)\quad 本位写8。
\item 十分位：\(\begin{aligned}[t]&\phantom{={}}2+1\\&=3\end{aligned}\)\quad 本位写3。
\item 个位：\(\begin{aligned}[t]&\phantom{={}}4+3\\&=7\end{aligned}\)\quad 本位写7。
\end{enumerate}
\par\smallskip
因此，\(\begin{aligned}[t]&\phantom{={}}4.2+3.18\\&=7.38\end{aligned}\)。
\par\medskip

\par\medskip\noindent
\textbf{第6小题：\(9.4-7.56=1.84\)}\par\smallskip
\noindent
\begin{tikzpicture}[x=1em,y=1em,baseline=(current bounding box.center),draw=answerred,text=answerred]
\node[anchor=base,inner sep=0pt] at (0.000,0) {$9$};
\node[anchor=base,inner sep=0pt] at (0.800,0) {$.$};
\node[anchor=base,inner sep=0pt] at (1.100,0) {$4$};
\node[anchor=base,inner sep=0pt] at (1.900,0) {$0$};
\node[anchor=base,inner sep=0pt] at (0.000,-1.1) {$7$};
\node[anchor=base,inner sep=0pt] at (0.800,-1.1) {$.$};
\node[anchor=base,inner sep=0pt] at (1.100,-1.1) {$5$};
\node[anchor=base,inner sep=0pt] at (1.900,-1.1) {$6$};
\node[anchor=base,inner sep=0pt] at (0.000,-2.3) {$1$};
\node[anchor=base,inner sep=0pt] at (0.800,-2.3) {$.$};
\node[anchor=base,inner sep=0pt] at (1.100,-2.3) {$8$};
\node[anchor=base,inner sep=0pt] at (1.900,-2.3) {$4$};
\node[anchor=base,inner sep=0pt] at (-.7,-1.1) {$-$};
\draw[line width=.5pt] (-1,-1.52)--(2.500,-1.52);
\draw[line width=.35pt] (-0.160,.03)--(0.160,.43);
\node[anchor=base,inner sep=0pt,font=\normalsize] at (0.000,.72) {$8$};
\draw[line width=.35pt] (0.940,.03)--(1.260,.43);
\node[anchor=base,inner sep=0pt,font=\normalsize] at (1.100,.72) {$13$};
\draw[line width=.35pt] (1.740,.03)--(2.060,.43);
\node[anchor=base,inner sep=0pt,font=\normalsize] at (1.900,.72) {$10$};
\end{tikzpicture}
\par\smallskip{\normalsize 划线数字是退位前的数，上方小字是退位后参与计算的数。\par}
\vspace{0pt}
\begin{enumerate}[leftmargin=1.7em,itemsep=5pt,topsep=0pt]
\item 十分位退位：\(\begin{aligned}[t]&\phantom{={}}4-1\\&=3\end{aligned}\)\quad 借出1给后一位。
\item 百分位接收退位：\(\begin{aligned}[t]&\phantom{={}}0+10\\&=10\end{aligned}\)\quad 前一位的1变成本位的10。
\item 百分位：\(\begin{aligned}[t]&\phantom{={}}10-6\\&=4\end{aligned}\)\quad 本位写4。
\item 个位退位：\(\begin{aligned}[t]&\phantom{={}}9-1\\&=8\end{aligned}\)\quad 借出1给后一位。
\item 十分位接收退位：\(\begin{aligned}[t]&\phantom{={}}3+10\\&=13\end{aligned}\)\quad 前一位的1变成本位的10。
\item 十分位：\(\begin{aligned}[t]&\phantom{={}}13-5\\&=8\end{aligned}\)\quad 本位写8。
\item 个位：\(\begin{aligned}[t]&\phantom{={}}8-7\\&=1\end{aligned}\)\quad 本位写1。
\end{enumerate}
\par\smallskip
因此，\(\begin{aligned}[t]&\phantom{={}}9.4-7.56\\&=1.84\end{aligned}\)。
\par\medskip

\par\medskip\noindent
\textbf{第6小题验算：\(1.84+7.56\)}\par\smallskip
\noindent
\begin{tikzpicture}[x=1em,y=1em,baseline=(current bounding box.center),draw=answerred,text=answerred]
\node[anchor=base,inner sep=0pt] at (0.000,0) {$1$};
\node[anchor=base,inner sep=0pt] at (0.800,0) {$.$};
\node[anchor=base,inner sep=0pt] at (1.100,0) {$8$};
\node[anchor=base,inner sep=0pt] at (1.900,0) {$4$};
\node[anchor=base,inner sep=0pt] at (0.000,-1.1) {$7$};
\node[anchor=base,inner sep=0pt] at (0.800,-1.1) {$.$};
\node[anchor=base,inner sep=0pt] at (1.100,-1.1) {$5$};
\node[anchor=base,inner sep=0pt] at (1.900,-1.1) {$6$};
\node[anchor=base,inner sep=0pt] at (0.000,-2.3) {$9$};
\node[anchor=base,inner sep=0pt] at (0.800,-2.3) {$.$};
\node[anchor=base,inner sep=0pt] at (1.100,-2.3) {$4$};
\node[anchor=base,inner sep=0pt] at (1.900,-2.3) {$0$};
\node[anchor=base,inner sep=0pt] at (-.7,-1.1) {$+$};
\draw[line width=.5pt] (-1,-1.52)--(2.500,-1.52);
\node[anchor=base,inner sep=0pt,font=\normalsize] at (1.100,.72) {$1$};
\node[anchor=base,inner sep=0pt,font=\normalsize] at (0.000,.72) {$1$};
\end{tikzpicture}
\par\smallskip{\normalsize 上方小字表示从后一位进来的数。\par}
\vspace{0pt}
\begin{enumerate}[leftmargin=1.7em,itemsep=5pt,topsep=0pt]
\item 百分位：\(\begin{aligned}[t]&\phantom{={}}4+6\\&=10\end{aligned}\)\quad 本位写0，向前一位进1。
\item 十分位：\(\begin{aligned}[t]&\phantom{={}}8+5+1\\&=13+1\\&=14\end{aligned}\)\quad 本位写4，向前一位进1。
\item 个位：\(\begin{aligned}[t]&\phantom{={}}1+7+1\\&=8+1\\&=9\end{aligned}\)\quad 本位写9。
\end{enumerate}
\par\smallskip
因此，\(\begin{aligned}[t]&\phantom{={}}1.84+7.56\\&=9.40\end{aligned}\)。
\par\medskip

\par\medskip\noindent
\textbf{第7小题：\(47.1-4.71=42.39\)}\par\smallskip
\noindent
\begin{tikzpicture}[x=1em,y=1em,baseline=(current bounding box.center),draw=answerred,text=answerred]
\node[anchor=base,inner sep=0pt] at (0.000,0) {$4$};
\node[anchor=base,inner sep=0pt] at (0.800,0) {$7$};
\node[anchor=base,inner sep=0pt] at (1.600,0) {$.$};
\node[anchor=base,inner sep=0pt] at (1.900,0) {$1$};
\node[anchor=base,inner sep=0pt] at (2.700,0) {$0$};
\node[anchor=base,inner sep=0pt] at (0.800,-1.1) {$4$};
\node[anchor=base,inner sep=0pt] at (1.600,-1.1) {$.$};
\node[anchor=base,inner sep=0pt] at (1.900,-1.1) {$7$};
\node[anchor=base,inner sep=0pt] at (2.700,-1.1) {$1$};
\node[anchor=base,inner sep=0pt] at (0.000,-2.3) {$4$};
\node[anchor=base,inner sep=0pt] at (0.800,-2.3) {$2$};
\node[anchor=base,inner sep=0pt] at (1.600,-2.3) {$.$};
\node[anchor=base,inner sep=0pt] at (1.900,-2.3) {$3$};
\node[anchor=base,inner sep=0pt] at (2.700,-2.3) {$9$};
\node[anchor=base,inner sep=0pt] at (-.7,-1.1) {$-$};
\draw[line width=.5pt] (-1,-1.52)--(3.300,-1.52);
\draw[line width=.35pt] (0.640,.03)--(0.960,.43);
\node[anchor=base,inner sep=0pt,font=\normalsize] at (0.800,.72) {$6$};
\draw[line width=.35pt] (1.740,.03)--(2.060,.43);
\node[anchor=base,inner sep=0pt,font=\normalsize] at (1.900,.72) {$10$};
\draw[line width=.35pt] (2.540,.03)--(2.860,.43);
\node[anchor=base,inner sep=0pt,font=\normalsize] at (2.700,.72) {$10$};
\end{tikzpicture}
\par\smallskip{\normalsize 划线数字是退位前的数，上方小字是退位后参与计算的数。\par}
\vspace{0pt}
\begin{enumerate}[leftmargin=1.7em,itemsep=5pt,topsep=0pt]
\item 十分位退位：\(\begin{aligned}[t]&\phantom{={}}1-1\\&=0\end{aligned}\)\quad 借出1给后一位。
\item 百分位接收退位：\(\begin{aligned}[t]&\phantom{={}}0+10\\&=10\end{aligned}\)\quad 前一位的1变成本位的10。
\item 百分位：\(\begin{aligned}[t]&\phantom{={}}10-1\\&=9\end{aligned}\)\quad 本位写9。
\item 个位退位：\(\begin{aligned}[t]&\phantom{={}}7-1\\&=6\end{aligned}\)\quad 借出1给后一位。
\item 十分位接收退位：\(\begin{aligned}[t]&\phantom{={}}0+10\\&=10\end{aligned}\)\quad 前一位的1变成本位的10。
\item 十分位：\(\begin{aligned}[t]&\phantom{={}}10-7\\&=3\end{aligned}\)\quad 本位写3。
\item 个位：\(\begin{aligned}[t]&\phantom{={}}6-4\\&=2\end{aligned}\)\quad 本位写2。
\item 十位：\(\begin{aligned}[t]&\phantom{={}}4-0\\&=4\end{aligned}\)\quad 本位写4。
\end{enumerate}
\par\smallskip
因此，\(\begin{aligned}[t]&\phantom{={}}47.1-4.71\\&=42.39\end{aligned}\)。
\par\medskip

\par\medskip\noindent
\textbf{第8小题：\(5.79+23.6=29.39\)}\par\smallskip
\noindent
\begin{tikzpicture}[x=1em,y=1em,baseline=(current bounding box.center),draw=answerred,text=answerred]
\node[anchor=base,inner sep=0pt] at (0.800,0) {$5$};
\node[anchor=base,inner sep=0pt] at (1.600,0) {$.$};
\node[anchor=base,inner sep=0pt] at (1.900,0) {$7$};
\node[anchor=base,inner sep=0pt] at (2.700,0) {$9$};
\node[anchor=base,inner sep=0pt] at (0.000,-1.1) {$2$};
\node[anchor=base,inner sep=0pt] at (0.800,-1.1) {$3$};
\node[anchor=base,inner sep=0pt] at (1.600,-1.1) {$.$};
\node[anchor=base,inner sep=0pt] at (1.900,-1.1) {$6$};
\node[anchor=base,inner sep=0pt] at (2.700,-1.1) {$0$};
\node[anchor=base,inner sep=0pt] at (0.000,-2.3) {$2$};
\node[anchor=base,inner sep=0pt] at (0.800,-2.3) {$9$};
\node[anchor=base,inner sep=0pt] at (1.600,-2.3) {$.$};
\node[anchor=base,inner sep=0pt] at (1.900,-2.3) {$3$};
\node[anchor=base,inner sep=0pt] at (2.700,-2.3) {$9$};
\node[anchor=base,inner sep=0pt] at (-.7,-1.1) {$+$};
\draw[line width=.5pt] (-1,-1.52)--(3.300,-1.52);
\node[anchor=base,inner sep=0pt,font=\normalsize] at (0.800,.72) {$1$};
\end{tikzpicture}
\par\smallskip{\normalsize 上方小字表示从后一位进来的数。\par}
\vspace{0pt}
\begin{enumerate}[leftmargin=1.7em,itemsep=5pt,topsep=0pt]
\item 百分位：\(\begin{aligned}[t]&\phantom{={}}9+0\\&=9\end{aligned}\)\quad 本位写9。
\item 十分位：\(\begin{aligned}[t]&\phantom{={}}7+6\\&=13\end{aligned}\)\quad 本位写3，向前一位进1。
\item 个位：\(\begin{aligned}[t]&\phantom{={}}5+3+1\\&=8+1\\&=9\end{aligned}\)\quad 本位写9。
\item 十位：\(\begin{aligned}[t]&\phantom{={}}0+2\\&=2\end{aligned}\)\quad 本位写2。
\end{enumerate}
\par\smallskip
因此，\(\begin{aligned}[t]&\phantom{={}}5.79+23.6\\&=29.39\end{aligned}\)。
\par\medskip

\par\medskip\noindent
\textbf{第9小题：\(30-9.29=20.71\)}\par\smallskip
\noindent
\begin{tikzpicture}[x=1em,y=1em,baseline=(current bounding box.center),draw=answerred,text=answerred]
\node[anchor=base,inner sep=0pt] at (0.000,0) {$3$};
\node[anchor=base,inner sep=0pt] at (0.800,0) {$0$};
\node[anchor=base,inner sep=0pt] at (1.600,0) {$.$};
\node[anchor=base,inner sep=0pt] at (1.900,0) {$0$};
\node[anchor=base,inner sep=0pt] at (2.700,0) {$0$};
\node[anchor=base,inner sep=0pt] at (0.800,-1.1) {$9$};
\node[anchor=base,inner sep=0pt] at (1.600,-1.1) {$.$};
\node[anchor=base,inner sep=0pt] at (1.900,-1.1) {$2$};
\node[anchor=base,inner sep=0pt] at (2.700,-1.1) {$9$};
\node[anchor=base,inner sep=0pt] at (0.000,-2.3) {$2$};
\node[anchor=base,inner sep=0pt] at (0.800,-2.3) {$0$};
\node[anchor=base,inner sep=0pt] at (1.600,-2.3) {$.$};
\node[anchor=base,inner sep=0pt] at (1.900,-2.3) {$7$};
\node[anchor=base,inner sep=0pt] at (2.700,-2.3) {$1$};
\node[anchor=base,inner sep=0pt] at (-.7,-1.1) {$-$};
\draw[line width=.5pt] (-1,-1.52)--(3.300,-1.52);
\draw[line width=.35pt] (-0.160,.03)--(0.160,.43);
\node[anchor=base,inner sep=0pt,font=\normalsize] at (0.000,.72) {$2$};
\draw[line width=.35pt] (0.640,.03)--(0.960,.43);
\node[anchor=base,inner sep=0pt,font=\normalsize] at (0.800,.72) {$9$};
\draw[line width=.35pt] (1.740,.03)--(2.060,.43);
\node[anchor=base,inner sep=0pt,font=\normalsize] at (1.900,.72) {$9$};
\draw[line width=.35pt] (2.540,.03)--(2.860,.43);
\node[anchor=base,inner sep=0pt,font=\normalsize] at (2.700,.72) {$10$};
\end{tikzpicture}
\par\smallskip{\normalsize 划线数字是退位前的数，上方小字是退位后参与计算的数。\par}
\vspace{0pt}
\begin{enumerate}[leftmargin=1.7em,itemsep=5pt,topsep=0pt]
\item 十位退位：\(\begin{aligned}[t]&\phantom{={}}3-1\\&=2\end{aligned}\)\quad 借出1给后一位。
\item 个位传递退位：\(\begin{aligned}[t]&\phantom{={}}0+10-1\\&=10-1\\&=9\end{aligned}\)\quad 收到10，再借出1给后一位。
\item 十分位传递退位：\(\begin{aligned}[t]&\phantom{={}}0+10-1\\&=10-1\\&=9\end{aligned}\)\quad 收到10，再借出1给后一位。
\item 百分位接收退位：\(\begin{aligned}[t]&\phantom{={}}0+10\\&=10\end{aligned}\)\quad 前一位的1变成本位的10。
\item 百分位：\(\begin{aligned}[t]&\phantom{={}}10-9\\&=1\end{aligned}\)\quad 本位写1。
\item 十分位：\(\begin{aligned}[t]&\phantom{={}}9-2\\&=7\end{aligned}\)\quad 本位写7。
\item 个位：\(\begin{aligned}[t]&\phantom{={}}9-9\\&=0\end{aligned}\)\quad 本位写0。
\item 十位：\(\begin{aligned}[t]&\phantom{={}}2-0\\&=2\end{aligned}\)\quad 本位写2。
\end{enumerate}
\par\smallskip
因此，\(\begin{aligned}[t]&\phantom{={}}30-9.29\\&=20.71\end{aligned}\)。
\par\medskip

\par\medskip\noindent
\textbf{第9小题验算：\(20.71+9.29\)}\par\smallskip
\noindent
\begin{tikzpicture}[x=1em,y=1em,baseline=(current bounding box.center),draw=answerred,text=answerred]
\node[anchor=base,inner sep=0pt] at (0.000,0) {$2$};
\node[anchor=base,inner sep=0pt] at (0.800,0) {$0$};
\node[anchor=base,inner sep=0pt] at (1.600,0) {$.$};
\node[anchor=base,inner sep=0pt] at (1.900,0) {$7$};
\node[anchor=base,inner sep=0pt] at (2.700,0) {$1$};
\node[anchor=base,inner sep=0pt] at (0.800,-1.1) {$9$};
\node[anchor=base,inner sep=0pt] at (1.600,-1.1) {$.$};
\node[anchor=base,inner sep=0pt] at (1.900,-1.1) {$2$};
\node[anchor=base,inner sep=0pt] at (2.700,-1.1) {$9$};
\node[anchor=base,inner sep=0pt] at (0.000,-2.3) {$3$};
\node[anchor=base,inner sep=0pt] at (0.800,-2.3) {$0$};
\node[anchor=base,inner sep=0pt] at (1.600,-2.3) {$.$};
\node[anchor=base,inner sep=0pt] at (1.900,-2.3) {$0$};
\node[anchor=base,inner sep=0pt] at (2.700,-2.3) {$0$};
\node[anchor=base,inner sep=0pt] at (-.7,-1.1) {$+$};
\draw[line width=.5pt] (-1,-1.52)--(3.300,-1.52);
\node[anchor=base,inner sep=0pt,font=\normalsize] at (1.900,.72) {$1$};
\node[anchor=base,inner sep=0pt,font=\normalsize] at (0.800,.72) {$1$};
\node[anchor=base,inner sep=0pt,font=\normalsize] at (0.000,.72) {$1$};
\end{tikzpicture}
\par\smallskip{\normalsize 上方小字表示从后一位进来的数。\par}
\vspace{0pt}
\begin{enumerate}[leftmargin=1.7em,itemsep=5pt,topsep=0pt]
\item 百分位：\(\begin{aligned}[t]&\phantom{={}}1+9\\&=10\end{aligned}\)\quad 本位写0，向前一位进1。
\item 十分位：\(\begin{aligned}[t]&\phantom{={}}7+2+1\\&=9+1\\&=10\end{aligned}\)\quad 本位写0，向前一位进1。
\item 个位：\(\begin{aligned}[t]&\phantom{={}}0+9+1\\&=9+1\\&=10\end{aligned}\)\quad 本位写0，向前一位进1。
\item 十位：\(\begin{aligned}[t]&\phantom{={}}2+0+1\\&=2+1\\&=3\end{aligned}\)\quad 本位写3。
\end{enumerate}
\par\smallskip
因此，\(\begin{aligned}[t]&\phantom{={}}20.71+9.29\\&=30.00\end{aligned}\)。
\par\medskip

}

\QHeading{第3题　解决问题。}

\Problem{小刚要买一个16.85元的文具盒，他带了20元，买完文具盒后还剩多少钱？}
\par


\MethodTitle{方法一}
\Answer{$20-16.85=3.15\text{（元）}$\par 答：还剩3.15元。}

\MethodTitle{方法二：用补差求剩余}
\Answer{从16.85元补到17元需0.15元，再从17元补到20元需3元。}
\Answer{\Formula{\begin{aligned}[t]&\phantom{={}}(17-16.85)+(20-17)\\
&=0.15+3\\
&=3.15\text{（元）}\end{aligned}}}
\Answer{答：还剩3.15元。}

